\documentclass[11pt]{article}
\usepackage{amsmath,amssymb}
\usepackage[margin=1in]{geometry}
\usepackage{hyperref}

\title{Ideas Parked --- Not Yet Merged (v14)}
\author{Staging document; each entry marks its eventual destination}
\date{}

\begin{document}
\maketitle

\noindent Three rough sketches, offered together as a single pass at
``what else might an atom's own internal structure predict,'' not as
three independent proposals. None is formalized here. Each is
checked against what the current paper already has --- one turns out
to restate a question already parked in v13; one turns out to be
partly already true of an existing metric, in a way worth stating
precisely rather than vaguely; one turns out to bear directly on an
already-open gap, as a specific, structurally motivated special case
of it. The sketches are kept in the user's own words where precision
would otherwise be lost, with the checking shown separately from the
sketch itself.

\section{Density, restated --- not a new entry}

\emph{``Density is number of total centers and info vectors
(possibly weighted by sum vector magnitude).''}

This is the same proposal as two conversations ago, in different
words: the formula candidates are the same ones already weighed in
ideas\_v13 (\(N\), \(D\), \(S_{\text{atom}}\), or
\(S_{\text{atom}}/N\)), and ``density'' is the same word already
retired once for \(\gamma_{\text{atom}}\) and recommended against for
any of these. Nothing here changes that weighing. Recorded only so a
future pass does not re-derive it: see ideas\_v13, Sections~3--5, for
the full accounting and the standing recommendation (keep \(N\)/\(D\)
as mass-number/atomic-number; leave ``density'' retired; park
\(S_{\text{atom}}/N\) rather than name it).

\section{Energy as shell distribution --- partly already true, worth
stating precisely}

\emph{``Energy is the distribution of strongest vectors\ldots in
upper shells vs lower shells. If an information atom has more
vectors in lower shells and less in higher shells and another has
more vectors in higher shells and less in lower shells then ones with
higher shells carry more energy.''}

\paragraph{Checked: this is already true of $\mathcal{E}_{\text{atom}}$,
for a reason worth stating rather than leaving implicit.}
$\mathcal{E}_{\text{atom}}=\sum r_i^2$ already rewards concentration
toward high $r$ over an even spread, for the same total strength ---
confirmed directly: four vectors summing to $S_{\text{atom}}=2.0$
give $\mathcal{E}_{\text{atom}}=1.0$ spread evenly
($0.5,0.5,0.5,0.5$), and $\mathcal{E}_{\text{atom}}=1.81$ concentrated
toward the extremes ($0.05,0.05,0.95,0.95$) --- an $81\%$ increase
for the identical total strength, driven entirely by where that
strength sits. This is not a coincidence of the example: squaring is
convex, so for any fixed sum, spreading values apart (in particular,
pushing some toward the high end) can only raise the sum of squares,
never lower it. $\sigma_{\text{atom}}$, checked on the same pair,
tracks the same shift ($0$ spread evenly, $0.333$ concentrated) ---
confirming $\mathcal{E}_{\text{atom}}$ and $\sigma_{\text{atom}}$ are
already telling a connected story, not two unrelated ones.

\paragraph{What is not yet true, and would need its own definition.}
$\sigma_{\text{atom}}$ is explicitly stated elsewhere in the paper to
be \emph{blind to which side} concentration sits on --- it measures
evenness versus unevenness, not high-shell versus low-shell. Checked
directly: nothing currently distinguishes ``three vectors low, one
high'' from ``three vectors high, one low'' when their evenness
happens to match, the way the sketch's own two-atom comparison
needs. $\mathcal{E}_{\text{atom}}$ partly closes this on its own, by
weight of the convexity argument above, but it is a single summed
number, not an explicit shell-indexed quantity; whether that is
enough, or whether a genuinely new, shell-position-weighted
quantity (something like a sum of $r_i^2$ each multiplied by its own
shell index) earns its own name under the ideas\_v13 rule, is not
decided here. A promising direction, not a finished one.

\section{Matter and antimatter --- a specific, structurally motivated
case of an already-open gap}

\emph{``If info atoms have all the vectors in right vertical sphere
then it is matter and another in left half then it is antimatter
particle with potential to cancel out all the information.''}

\paragraph{Checked: this bears directly on Gap~1, not on new ground.}
The paper already has one proven cancellation: a mark and an
equal-strength anti-phase mark, same direction $\hat n$, Bundle as
two separate seats ($\uplus$ does not cancel them) and are only
cancelled on a later reading, by Mix, into black (``Pathogen
cancellation,'' a costume of Axioms~4 and~5). That is opposition in
$\chi$ at one shared direction --- a narrow, well-defined special
case, not a general rule for differing directions.

The sketch proposes a second, geometrically distinct kind of
opposition: not same-$\hat n$, opposite-$\chi$, but opposite-$\hat n$
itself --- antipodal points on the sphere ($\hat n$ against $-\hat
n$). This is exactly the question ``Open: Mix across differing
directions'' (Gap~1) leaves unsettled, and the leading candidate on
file for that gap, Route~1, argues for refusing combination across
differing directions \emph{in general}, on evidence from the
$\mathtt{chladni}$ inclusion. The sketch is not in conflict with
Route~1 so much as it is the same kind of move the existing
cancellation rule already made: Route~1 is about the generic case
(some unprivileged angle apart), while both the existing
same-direction cancellation and this sketch's antipodal cancellation
are specific, maximally symmetric cases --- $0^\circ$ apart in $\chi$
reversed, and $180^\circ$ apart in $\hat n$ exactly --- each singular
enough on the sphere to plausibly warrant its own rule rather than
inheriting the generic refusal.

\paragraph{What would be needed before this is more than a sketch.}
A stated rule for what Mix (or some other named reading) does at
exactly $\hat n_B=-\hat n_A$, checked the way the existing
cancellation was checked --- does it reduce to black the way
same-direction opposite-phase does, or to something else; does it
require equal strength, the way the existing case implicitly does
(an equal-strength anti-phase mark); and whether ``cancels out all
the information'' means the pair reduces to the Null State of
Axiom~2 specifically, or some other outcome not yet named. None of
this is decided here --- the sketch is recorded as a specific,
well-motivated candidate for Gap~1's hardest symmetric case, not as
a solved instance of it.

\section{Status}

All three sketches are the user's own framing, offered as a rough
pass with the explicit caveat that it is not yet known which will
hold up. Section~1 is closed by reference to prior work, not by new
reasoning. Sections~2 and~3 are each partly verified (the numbers
and the citations above are checked against the current paper) and
partly open (the shell-weighted formula in Section~2; the exact
cancellation rule in Section~3), and both are left open rather than
pushed toward a formalization the sketches themselves do not yet
claim to have earned.

\end{document}
